% =====================================================================
%  THÈME 4 — Logarithmes, pH & incertitude — CORRIGÉ — Niveau DIFFICILE
% =====================================================================
\documentclass[11pt,a4paper]{article}
\usepackage[utf8]{inputenc}
\usepackage[T1]{fontenc}
\usepackage{lmodern}
\usepackage[french]{babel}
\usepackage{amsmath,amssymb,mathtools}
\usepackage{icomma}
\usepackage{siunitx}
\sisetup{output-decimal-marker={,},per-mode=symbol}
\usepackage[version=4]{mhchem}
\usepackage{xcolor}
\usepackage[most]{tcolorbox}
\usepackage{enumitem}
\usepackage{fancyhdr}
\usepackage[a4paper,margin=2.2cm,headheight=26pt,headsep=8pt]{geometry}

\definecolor{primary}{RGB}{150,98,162}
\definecolor{primarydark}{RGB}{92,52,110}
\newcommand{\cs}{chiffre significatif}
\newcommand{\CS}{chiffres significatifs}

\pagestyle{fancy}\fancyhf{}
\fancyhead[L]{\small\textcolor{primarydark}{\textbf{Fiche 1} — Chiffres significatifs}}
\fancyhead[R]{\small\textcolor{primarydark}{\textsc{Corrigé · Difficile · Thème 4}}}
\fancyfoot[C]{\small\thepage}
\renewcommand{\headrulewidth}{0.8pt}
\renewcommand{\headrule}{\hbox to\headwidth{\color{primary}\leaders\hrule height \headrulewidth\hfill}}

\newtcolorbox[auto counter]{cor}[1][]{enhanced,breakable,boxrule=0.8pt,arc=2pt,
  colback=primarydark!5,colframe=primarydark,left=3mm,right=3mm,top=2mm,bottom=2mm,
  fonttitle=\bfseries,coltitle=white,
  title={Correction — Exercice~\thetcbcounter\ifx\relax#1\relax\relax\else\ ---\ \textmd{\itshape #1}\fi}}

\setlength{\parindent}{0pt}\setlength{\parskip}{3pt}

\begin{document}

\begin{center}
{\Large\bfseries\color{primarydark} Corrigé — Niveau Difficile}\\[1pt]
{\color{primary} Thème 4 : pH \& incertitude — synthèse}
\end{center}
\vspace{2mm}

\begin{cor}[pH dans les deux sens, et chiffres réellement significatifs]
\begin{enumerate}[label=\alph*),leftmargin=1.6em,itemsep=1pt]
  \item $\mathrm{pH}=-\log(\num{3,0e-5})=5-\log(\num{3,0})=5-\num{0,48}=\num{4,52}$. La concentration
        a $2$~CS $\Rightarrow$ pH à \textbf{2 décimales} : $\boxed{\mathrm{pH}=\num{4,52}}$.
  \item Les chiffres significatifs du pH sont \emph{les deux décimales} « $\mathbf{52}$ » ($2$ chiffres,
        comme les $2$~CS de la concentration). La partie entière « $4$ » n'est \emph{pas} un CS : elle
        ne fait que rappeler l'ordre de grandeur ($10^{-5}$).
  \item Vérification : $[\ce{H3O+}]=10^{-4,52}=10^{\num{0,48}}\times10^{-5}=\num{3,0e-5}\,\si{\mol\per\litre}$
        ($2$~CS). On retombe bien sur la donnée initiale.
\end{enumerate}
\end{cor}

\begin{cor}[Deux balances : beaucoup de CS ne veut pas dire « juste »]
\begin{enumerate}[label=\alph*),leftmargin=1.6em,itemsep=1pt]
  \item Balance A : $\Delta=\qty{0,0001}{\gram}$, $\dfrac{\Delta}{x}=\dfrac{\num{0,0001}}{\num{2,0500}}
        \approx\qty{0,005}{\percent}$. \\
        Balance B : $\Delta=\qty{0,01}{\gram}$, $\dfrac{\Delta}{x}=\dfrac{\num{0,01}}{\num{2,00}}
        =\qty{0,5}{\percent}$.
  \item La plus \textbf{fidèle} est \textbf{A} ($5$~CS, incertitude relative $100$ fois plus faible).
  \item La plus \textbf{juste} est \textbf{B} : son écart à la vraie valeur est $\le\qty{0,01}{\gram}$
        ($\num{2,00}$ contient $\num{2,000}$), alors que A affiche $\num{2,0500}$, soit un biais de
        $+\qty{0,05}{\gram}$ ($=\qty{50}{\milli\gram}$) bien au-delà de sa propre incertitude.
  \item \textbf{Non.} Un grand nombre de \CS\ traduit la \emph{fidélité}, pas la \emph{justesse}. Une
        balance déréglée reste fausse quel que soit le nombre de chiffres qu'elle affiche : les CS ne
        corrigent jamais un biais.
\end{enumerate}
\end{cor}

\begin{cor}[Un titrage : de l'incertitude au pH]
\begin{enumerate}[label=\alph*),leftmargin=1.6em,itemsep=1pt]
  \item $\dfrac{\Delta[\ce{H3O+}]}{[\ce{H3O+}]}=\dfrac{\num{0,1}}{\num{2,0}}=\num{0,05}=\qty{5}{\percent}$.
        Une incertitude relative de l'ordre du pourcent correspond à $\mathbf{2}$~CS : on écrit donc la
        concentration avec $2$~CS ($\num{2,0e-3}$), et surtout pas davantage.
  \item $\mathrm{pH}=-\log(\num{2,0e-3})=3-\num{0,30}=\num{2,70}$ ; $2$~CS $\to 2$ décimales
        $\Rightarrow \boxed{\mathrm{pH}=\num{2,70}}$.
  \item Dans un logarithme, seule la partie décimale porte l'information significative (la partie
        entière code la puissance de dix) : le nombre de \emph{décimales} du pH égale donc le nombre
        de \emph{\CS} de $[\ce{H3O+}]$.
\end{enumerate}
\end{cor}

\end{document}
